\documentclass{article}
\usepackage{amsmath}
\usepackage{amssymb}
\usepackage{amsthm}
\newtheorem{theorem}{Theorem}[section]
\newtheorem{lemma}[theorem]{Lemma}
\renewcommand*{\proofname}{Proof}
\theoremstyle{definition}
\newtheorem{definition}{Definition}[section]
\theoremstyle{remark}
\newtheorem*{remark}{Remark}
\theoremstyle{example}
\newtheorem{example}{Example}[section]
\DeclareMathOperator*{\argmax}{arg\,max}
\DeclareMathOperator*{\argmin}{arg\,min}
\usepackage{fullpage}
%\usepackage{svg}
\usepackage{pgf}
\usepackage{mathtools}
\usepackage{physics}
\usepackage{outlines}
\allowdisplaybreaks

\newcommand{\AAA}{\mathbf{A}}
\newcommand{\uu}{\mathbf{u}}
\newcommand{\vv}{\mathbf{v}}
\newcommand{\betat}{\hat{\beta}}
\newcommand{\mx}{\overline{x}}
\newcommand{\my}{\overline{y}}
\newcommand{\mxy}{\overline{x}\overline{y}}
\newcommand{\mxsquare}{\overline{x^2}}
\newcommand{\e}{\mathbf{e}}
\newcommand{\x}{\mathbf{x}}

\newcommand{\lambdafunction}[2][1]{
		\let\tmp\relax
		\newcommand{\tmp}[#1]{#2}
		\tmp
}

\newcommand{\ffmatrix}[3]{
		\begin{bmatrix}
				\lambdafunction[2]{#1}{1}{1} & \lambdafunction[2]{#1}{1}{2} & \cdots & \lambdafunction[2]{#1}{1}{#3}\\
				\lambdafunction[2]{#1}{2}{1} & \lambdafunction[2]{#1}{2}{2} & \cdots & \lambdafunction[2]{#1}{2}{#3}\\
				\vdots & \vdots & \ddots & \vdots\\
				\lambdafunction[2]{#1}{#2}{1} & \lambdafunction[2]{#1}{#2}{2} & \cdots & \lambdafunction[2]{#1}{#2}{#3}\\
		\end{bmatrix}
}

\newcommand{\fmatrix}[3]{\begin{bmatrix} #1_{1 1} & #1_{1 2} & \cdots & #1_{1 #3} \\ #1_{2 1} & #1_{2 2} & \cdots & #1_{2 #3} \\ \vdots & \vdots & \ddots & \vdots \\ #1_{#2 1} & #1_{#2 2} & \cdots & #1_{#2 #3} \end{bmatrix}}
\newcommand{\fvector}[2]{\begin{bmatrix} #1_1 \\ #1_2 \\ \vdots \\ #1_{#2} \end{bmatrix}}

\newcommand{\mycomment}[1]{}

\let\P\relax
\NewDocumentCommand{\P}{o}{\mathbb{P}\IfValueT{#1}{\left[#1\right]}}
\newcommand{\E}[1]{\mathbb{E}[#1]}



\title{Learning Theory}
\author{Lysandre Terrisse}

\begin{document}
\maketitle

\section{More formally}

%Thanks to Robust Machine Learning, chapter 2.
%Also thanks to Mathematical Aspects of Deep Learning
%Also thanks to Learning Theory: An Approximation Theory Viewpoint

%\begin{definition}[Supervised problem]
%		Let $\mathcal{X}$ and $\mathcal{Y}$ be sets, respectively called the input and output spaces, and let $\mathcal{D}$ be a probability distribution over $\mathcal{X} \times \mathcal{Y}$. The goal of a supervised problem is, given a finite \textit{training set} $\mathcal{S} \in (\mathcal{X} \times \mathcal{Y})^m$ of size $m$, to find a model $f : \mathcal{X} \rightarrow \mathcal{Y}$ which minimizes the risk $\mathcal{R}(f)$ for a given \textit{loss function} $\mathcal{L} : (\mathcal{X} \rightarrow \mathcal{Y}) \times (\mathcal{X} \times \mathcal{Y}) \rightarrow \mathbb{R}$:
%		$$\mathcal{R}(f) \coloneq \mathbb{E}_{(x, y) \sim \mathcal{D}}[\mathcal{L}(f, (x, y))] = \int_{\mathcal{X} \times \mathcal{Y}} \mathcal{L}(f, (x, y)) \dd \mathbb{P}$$
%\end{definition}

\begin{definition}[Supervised problem]
		Let $(\mathcal{X} \times \mathcal{Y}, \mathcal{A}, \P)$ be a measurable space, with $\mathcal{X}$ and $\mathcal{Y}$ respectively called the \textit{input space} and \textit{output space}. The goal of a supervised problem is, given a finite \textit{training set} $\mathcal{S} \in (\mathcal{X} \times \mathcal{Y})^m$ of size $m$, and a subset $\mathcal{F}$ of $\mathcal{X} \rightarrow \mathcal{Y}$ called the \textit{hypothesis set}, to find a \textit{model} $f \in \mathcal{F}$ which minimizes the risk $\mathcal{R}(f)$ for a given integrable\footnote{For $\E{\mathcal{L}(f, (x, y))} = \int_{\mathcal{X} \times \mathcal{Y}} \mathcal{L}(f, (x, y)) \dd \mathbb{P}$ to always be defined, it is necessary and sufficient that $\mathcal{L}(f, (x, y))$ is integrable on its second argument. For many loss functions, such as the quadratic loss function, this condition is equivalent to stating that each function $f$ belongs to $L^2(\mathcal{X} \times \mathcal{Y})$ and that $\E{y^2} = \int_{\mathcal{X} \times \mathcal{Y}} y^2 \dd \P < \infty$. These conditions in turn imply for the Bayes optimal function that $\E{Y \mid X} \in L^2(\mathcal{X} \times \mathcal{Y})$.} \textit{loss function} $\mathcal{L} : (\mathcal{X} \rightarrow \mathcal{Y}) \times (\mathcal{X} \times \mathcal{Y}) \rightarrow \mathbb{R}$:
		$$\mathcal{R}(f) \coloneq \E{\mathcal{L}(f, (x, y))}$$
\end{definition}


\begin{definition}[Supervised regression problem]
		A supervised \textit{regression problem} is a supervised problem where $\mathcal{Y} \subseteq \mathbb{R}$, where $\mathcal{F} \subseteq L^2(\mathcal{X} \times \mathcal{Y})$, where $\E{y^2} < \infty$, and where the loss function is the \textit{quadratic loss function}:
				$$\mathcal{L}(f, (x, y)) = (f(x) - y)^2$$
\end{definition}
%For the quadratic loss function to always be integrable, we also need each $f$ to be integrable.

%Est-ce qu'il faut admettre que E[(f(x) - f*(x))²] et E[(f²(x) - y)²] sont finis
%Quelles choses doivent être mesurables (coût, modèles)
%Comment montrer que l'intégrale des P[y | x] selon y vaut 1
\begin{theorem}[Bayes optimal function]
		If we knew the probability distribution $\P$, an optimal solution of the supervised regression problem would be the \textit{Bayes optimal function} $f^*$, defined arbitrarily for each $x \in \mathcal{X}$ such that $\P[x] = 0$, and otherwise defined as:
		$$f^*(x) \coloneq \E{y \mid x}$$% = \int_\mathcal{Y} y \dd \P[y \mid x]$$
		%with the following definitions:
		%\begin{align*}
		%		\mathbb{E}_{(x, y) \sim \mathcal{D}}[y \mid x] &\coloneq \int_\mathcal{Y} y \mathbb{P}_\mathcal{D}[y \mid x] \dd y &
		%		\mathbb{P}_\mathcal{D}[y \mid x] &\coloneq \frac{\mathbb{P}_\mathcal{D}[x, y]}{\mathbb{P}_\mathcal{D}[x]} &
		%		\mathbb{P}_\mathcal{D}[x] &\coloneq \int_\mathcal{Y} \mathbb{P}_\mathcal{D}[x, y] \dd y
		%\end{align*}
		To see that the Bayes optimal solution indeed is a minimum of the risk function, we will prove that, for any function $f : \mathcal{X} \rightarrow \mathcal{Y}$, we can decompose the risk like so:
				$$\mathcal{R}(f) = \E{(f(x) - f^*(x))^2} + \mathcal{R}(f^*)$$
		which indeed reaches a minimum when $f = f^*$.
\end{theorem}

\begin{proof}
		Firstly, let's prove that for all $x \in \mathcal{X}$, we have that $\int_\mathcal{Y} (f^*(x) - y)\P[x, y] \dd y = 0$. If $\P[x]=0$, then the equality is trivially verified. If $\P[x] \neq 0$, we have:
		\begin{align*}
				\int_\mathcal{Y} (f^*(x) - y)\P[x, y] \dd y &= \int_\mathcal{Y} (f^*(x) - y)\P[y \mid x] \P[x] \dd y\\
																				&= \P[x] \int_\mathcal{Y} (f^*(x) - y)\P[y \mid x] \dd y\\
																				&= \P[x] \left(\int_\mathcal{Y} f^*(x)\P[y \mid x] \dd y - \int_\mathcal{Y} y\P[y \mid x] \dd y\right)\\
																				&= \P[x] \left(f^*(x) \int_\mathcal{Y} \P[y \mid x] \dd y - \int_\mathcal{Y} y\P[y \mid x] \dd y\right)\\
																				&= \P[x] \left(f^*(x) - \int_\mathcal{Y} y\P[y \mid x] \dd y\right)\\
																				&= \P[x] \cdot 0\\
																				&= 0
		\end{align*}
		The risk of the supervised regression problem can now be rewritten as:
		\begin{align*}
				\mathcal{R}(f) &= \E{(f(x) - y)^2}\\
							   &= \E{(f(x) - f^*(x) + f^*(x) - y)^2}\\
							   &= \E{(f(x) - f^*(x))^2 + (f^*(x) - y)^2 + 2(f(x) - f^*(x))(f^*(x) - y)}\\
							   &= \E{(f(x) - f^*(x))^2} + \mathcal{R}(f^*) + \int_\mathcal{X} \int_\mathcal{Y} 2(f(x) - f^*(x))(f^*(x) - y) \P[x, y] \dd y \dd x\\
							   &= \E{(f(x) - f^*(x))^2} + \mathcal{R}(f^*) + \int_\mathcal{X} \left(2(f(x) - f^*(x)) \int_\mathcal{Y} (f^*(x) - y) \P[x, y] \dd y \right) \dd x\\
							   &= \E{(f(x) - f^*(x))^2} + \mathcal{R}(f^*) + \int_\mathcal{X} \left(2(f(x) - f^*(x)) \cdot 0 \right) \dd x\\
							   &= \E{(f(x) - f^*(x))^2} + \mathcal{R}(f^*)
		\end{align*}		
\end{proof}

\begin{definition}[Empirical risk minimization]
		In practice, we do not know the distribution $\P$, and we instead just have access to the training set $\mathcal{S}$. In that case, one strategy for choosing the model $f$ is the \textit{empirical risk minimization}, which consists of choosing the model which minimizes the empirical risk $\widehat{\mathcal{R}}_\mathcal{S}(f)$:
				$$\widehat{\mathcal{R}}_\mathcal{S}(f) \coloneq \frac{1}{m} \sum_{(x, y) \in \mathcal{S}} \mathcal{L}(f, (x, y))$$
\end{definition}

%\begin{definition}[Supervised binary classification problem]
%		A supervised \textit{binary classification problem} is a supervised problem where $|\mathcal{Y}|=2$ and where the loss function is the \textit{0-1 loss function}:
%				$$f, (x, y) \mapsto \begin{cases} 0 & \tif f(x) = y \\ 1 & \tif f(x) \neq y \end{cases}$$
%		In that case, the risk is equal to:
%				$$\mathcal{R}(f) = \mathbb{P}_{(x, y) \sim \mathcal{D}}[f(x) \neq y]$$
%\end{definition}

%\begin{theorem}[Bayes optimal classifier]
%		If we knew the distribution $\mathcal{D}$, the optimal solution of supervised binary classification problems when using the misclassification error would be the \textit{Bayes optimal classifier} $f^*$:
%				$$f^*(x) = \argmax_{y \in \mathcal{Y}} \mathbb{P}_{(x, y) \sim \mathcal{D}}[y \mid x]$$
%\end{theorem}



\end{document}
